\documentclass{article}
\usepackage[utf8]{inputenc}
\usepackage{graphicx}
\usepackage[a4paper,top=2.5cm,bottom=2cm,left=1.5cm,right=1.5cm,marginparwidth=1.75cm]{geometry}
\usepackage{tikz}
\usetikzlibrary{shapes.misc,shadows}



\begin{document}
\begin{titlepage}
\begin{center}
    \includegraphics[height=4cm,width=4cm]{NSF.png}\\
    \vspace{10mm} %5mm vertical space
    \huge{NSF Project Report Format}\\
    \vspace{10mm} %5mm vertical space
    \huge{Principal Investigator}\\
    \vspace{10mm} %5mm vertical space
    \large{\today}
\end{center}
\end{titlepage}
\newpage
%\tableofcontents
%\listoffigures 
\newpage
\section*{Accomplishments}
%You have the option of selecting “nothing to report” in this section.
\subsection*{What are the major goals of the project?}
\vspace{5mm}
\subsection*{What was accomplished under these goals (you must provide information for at least one of the 4 categories below)?}
\vspace{5mm}
\begin{itemize} 
\setlength{\itemsep}{5pt}
 \setlength{\parskip}{5pt}
 \setlength{\parsep}{5pt}
    \item[ ] Major Activities:\\
    \item[ ] Specific Objectives:\\
    \item[ ] Significant Results:\\
    \item[ ] Key outcomes or other achievements:\\
\end{itemize}
\subsection*{What opportunities for training and professional development has the project provided?}
\vspace{5mm}
\subsection*{How have the results been disseminated to communities of interest? }
\vspace{5mm}
\subsection*{What do you plan to do during the next reporting period to accomplish the goals?}
\vspace{5mm}
\newpage
\section*{Products}
%You have the option of selecting “nothing to report” in this section. There are no limitations to the number of entries you submit and you can also pull information directly from Thomson Search when using the online tool on Research.gov.
%Within the Products section, you can list any products resulting from your project during the specified reporting period, such as:
\begin{itemize}
\setlength{\itemsep}{5pt}
 \setlength{\parskip}{5pt}
 \setlength{\parsep}{5pt}
    \item[ ] Journals or Juried Conference Papers
    \item[ ] Books:
    \item[ ] Thesis/Dissertations
    \item[ ] Other Conference Presentations / Papers
    \item[ ] Other Publications
    \item[ ] Technologies or Techniques
    \item[ ] Patents
    \item[ ] Inventions
    \item[ ] Licenses
    \item[ ] Websites
    \item[ ] Other Products
\end{itemize}
%NOTE: You may upload PDF files with images, tables, charts, or other graphics in support of the Products section. You may upload up to 4 PDF files with a maximum file size of 5 MB each
\newpage
\section*{Participants}
%There are no limits on the number of participants you list for this section; however, you must list participants who have worked one person month or more for the project reporting period. You have the option of selecting “nothing to report” in this section. For Research Experience for Undergraduates (REU) sites and supplements, specific questions will be listed in this section. The online service will also ask for additional information on participants such as: 
%What individuals have worked on the project?
%What organizations have been involved as partners


\subsection*{What individuals have worked on the project?}
\begin{tabular}{c@{\hskip 1.5in} c@{\hskip 0.25in} c} 
Name & Most Senior Project Role & Nearest Person Month Worked 
\\\hline
\\\hline
\\\hline
\end{tabular}
\vspace{5mm}
\subsection*{What other organizations have been involved as partners?}
%The online service will also ask you for additional information such as:

\begin{itemize}
\setlength{\itemsep}{5pt}
 \setlength{\parskip}{5pt}
 \setlength{\parsep}{5pt}
    \item Type of Partner Organization
    \item Name
    \item Location
    \item Partner’s contribution to the project
\end{itemize}
\vspace{5mm}
\subsection*{Have other collaborators or contacts been involved?} 
\begin{tikzpicture}[baseline=(char.base)]
\node(char)[draw,fill=white,
  shape=rounded rectangle,
  drop shadow={opacity=.5,shadow xshift=0pt},
  minimum width=1.8cm]
  {Yes};
\end{tikzpicture} \hspace{1in} No
%Yes \hspace{1in}
%\begin{tikzpicture}[baseline=(char.base)]
%\node(char)[draw,fill=white,
%  shape=rounded rectangle,
%  drop shadow={opacity=.5,shadow xshift=0pt},
%  minimum width=1.8cm]
%  {No};
%\end{tikzpicture}
\newpage
\section*{Impacts}
%You have the option of selecting “nothing to report” in this section.
\subsection*{What is the impact on the development of the principal discipline(s) of the project?}
\vspace{5mm}
\subsection*{What is the impact on other disciplines?}
\vspace{5mm}
\subsection*{What is the impact on the development of human resources?}
\vspace{5mm}
\subsection*{What is the impact on institutional resources that form infrastructure?}
\vspace{5mm}
\subsection*{What is the impact on information resources that form infrastructure?}
\vspace{5mm}
\subsection*{What is the impact on technology transfer?}
\vspace{5mm}
\subsection*{What is the impact on society beyond science and technology?}
\newpage
\section*{Changes / Problems}
%If not previously reported in writing to the agency through other mechanisms, provide the following additional information or state, "Nothing to Report", if applicable.

\subsection*{Changes in approach and reason for change:}
\vspace{5mm}
\subsection*{Actual or Anticipated problems or delays and actions or plans to resolve them:}
\vspace{5mm}
\subsection*{Changes that have a significant impact on expenditures:}
\vspace{5mm}
\subsection*{Significant changes in use or care of human subjects:}
\vspace{5mm}
\subsection*{Significant changes in use or care of biohazards:}
\vspace{5mm}
\newpage
\section*{Special Requirements}
%This report section is only available when Special Requirements are specifically noted in the solicitation and approved by the Office of Management and Budget.
%NOTE: You may upload PDF files in support of the Special Requirements section. You may upload PDF files with a maximum file size of 10 MB each. There is no limit to the number of files uploaded. 
\end{document}